% Pursuing Stacks (cotes 134-2 to 134-8): archive cote 134-7, site page 11 (PDF page 12), typescript page 453.
% typescript. Draft compared: ps6.tex:491-560.
% AI-transcribed by Claude Opus 5.5 (claude-opus-5-5, adaptive thinking, effort high), Message Batches API, batch ps3 (msgbatch_01XdFmuVVBAcgc7EFxdtCqsj), 2026-11-12; not yet checked by eye.
\scanpage{134-7/11}
\mspage{453}

I am not too sure whether all this is'nt just complete nonsense -
it is worth getting it clear whether it is or not, at any rate ! There
is one case of special interest, the ``simplest'' one in a way, namely
when \xout{Xx} the simplicial maps between the $X_n$'s , each represented by a
k-module $M_n$ , are in fact k-linear, in other words, when $X_{\hand{\ast}}$ comes
from a semisimplicial k-module $M_{\hand{\ast}}$ - more specifically still, when
this is a $K(\hand{\pi},n)$ type, say the nicest semisimplicial model of this,
using the Kan-Dold-Puppe functor for the chain complex of k-modules,
having $\hand{\pi}$ in degree n and zero elsewhere. Then the left hand side of u)
gives rise to the Eilenberg-Maclane homology groups
\[ \hand{(v)} \qquad H_i(\hand{\pi},n;\underline{Z}) \quad , \]
\margin{$\times$} which I guess should be k-modules a priori, because of the operations
of k upon $\hand{\pi}$ , and the question then arises whether these can be compu-%
ted using the right hand complex $L(X_{\hand{\ast}})$ . Maybe such a thing is even a
familiar fact for people in the know ? If it turned out to be false
even for $k=\underline{Z}$ , my faith in schematization of homotopy types would be
seriously shaken I confess ...

\hand{* * *}

After a little break for dinner, just one more afterthought. Wor-%
king with the completions $\Gamma_k\hat{}$ may seem a little forbidding, and all
the more so if it is used for computing \uline{homology} invariants, not cohomo-%
logy (the latter more likely to involve \inserted[typed]{infinite} products in \xout{\unsure{the}} the correspon-%
ding (cochain) complexes ...). On the other hand, as was explicitly
stated from the beginning, the natural context here seems to be
\uline{pointed} homotopy types, and hence ``pointed'' algebraic para\xout{\unsure{g}}digms for
these - an aspect which we lost sight of, when looking for a suitable
description of some category U(k) of ``unipotent bundles'' over k. It
would seem that in the ``question'' above, we should therefore insist
that $X_{\hand{\ast}}$ should be a semi-simplicial object of the category $U(k)^{\bullet}$ of
``pointed'' objects of U(k), namely objects \inserted{X} endowed with a section over
the final object e (i.e. with an \deleted{object of} \inserted{element in} X(k)). This will be automa-%
tic at any rate in terms of the condition \hand{(}t\hand{)} , $X_o = e$ . The point I
wish to make is that the category of pointed objects of U(k) admits a
somewhat simpler description (by cho\inserted{o}sing the marked point as the ``origin''
for \deleted{a given} parametrisation of \unsure{\hand{the}} \inserted{given} object X of U(k) by a k-module M),
by model-objects which are still arbitrary k-modules M, but the ``maps''
now being k-\uline{linear} \deleted{maps} \inserted{\uline{continuous maps}}
\[ \hand{(}p') \qquad f : \Gamma_k(M) \hand{\longrightarrow} \Gamma_k(M') \]

\margin{11}
